\subsection{The second-order state ceiling is attained}
The order $B^{-2}$ can be attained by resource programs of the same running-mean learner.
Take a single Gaussian regime with $\theta=1/3$, $\lambda=1/5$, $Q=8$, $h=8$, $D=0$ and $r^2=49/50$.
A two-observation all-actor probe uses $32$ messages and two clock units.
With total budgets $(80B+32,32B+2)$, its two continuing programs $(q,b)=(2,1)$ and $(8,1)$ have applied horizons $8B$ and $5B$.
Both have leading risk $1/(8B)$.
The invariant contrast is one dimensional, with explained terminal-state variance
\[
s^2=\frac{[\theta(1-\lambda)]^2}{2[\theta(1-\lambda)+1/8]}=\frac{64}{705}.
\]
Let $c$ select the smaller unconditional commitment risk and $u$ the even conditional-risk minimum over these two programs.
For all sufficiently large integer $B$, $c$ selects $(2,1)$.
The exact geometric risk formula yields
\[
B^2(J_{8,1}^{r}-J_{2,1}^{r})\longrightarrow a=\frac{19}{320000},
\]
whereas $B^2(k_{8,1}^2-k_{2,1}^2)s^2\longrightarrow b=13/94000$.
Here $J_{q,b}^{r}$ and $k_{q,b}$ name commitment risk and state transfer, not a new model index.
Since $0<a<b$, the conditional scores cross on a set of positive Gaussian probability.
Writing $G\sim\mathcal N(0,1)$, pointwise minimization and dominated convergence give
\[
\lim_{B\to\infty}B^2[R(c)-R(u)]
=b\,\mathbb E\left[\left(\frac{1187}{2080}-G^2\right)_+\right]>0.
\]
For domination, the scaled risk gaps and state coefficients are bounded by the geometric expansions, so the positive part is bounded by a constant times $1+G^2$.
Thus the state-value ceiling has the sharp budget order in this program class, even with no unknown regime to identify.
The construction establishes strict second-order value, not a first-order gain or confirmation of the separate empirical state-selection test.
